% Part: incompleteness
% Chapter: representability-in-q
% Section: representable-comp

\documentclass[../../../include/open-logic-section]{subfiles}

\begin{document}

\olfileid[hi]{inc}{req}{rpc}
\olsection{$\Th{Q}$ में निरूपणीय फलन संगणनीय हैं}

हम सिद्ध करेंगे कि~$\Th{Q}$ में निरूपणीय प्रत्येक फलन संगणनीय है। पहले
हमें~$\Th{Q}$ में निरूपणीय फलनों के बारे में एक उपपत्ति स्थापित करनी होगी।

\begin{lem}\ollabel{lem:rep-q}
  यदि $f(x_0, \dots, x_k)$,~$\Th{Q}$ में निरूपणीय है, तो कोई
  !!a{formula}~$!A(x_0, \dots, x_k, y)$ ऐसा है कि
  \[
  \Th{Q} \Proves !A_f(\num{n_0}, \dots, \num{n_k}, \num{m})
  \quad\text{तभी और केवल तभी जब}\quad m = f(n_0, \dots, n_k).
  \]
\end{lem}

\begin{proof}
  ``यदि'' भाग
  \olref[int]{defn:representable-fn}\olref[int]{defn:rep:a} है। ``केवल
  यदि'' भाग इस प्रकार मिलता है: मान लें कि
  $\Th{Q} \Proves !A_f(\num{n_0}, \dots, \num{n_k}, \num{m})$, किंतु
  $m \neq f(n_0, \dots, n_k)$। मानें $l = f(n_0, \dots, n_k)$।
  \olref[int]{defn:representable-fn}\olref[int]{defn:rep:a} के अनुसार
  $\Th{Q} \Proves !A_f(\num{n_0}, \dots, \num{n_k}, \num{l})$।
  \olref[int]{defn:representable-fn}\olref[int]{defn:rep:b} के अनुसार
  $\lforall[y][(!A_f(\num{n_0}, \dots, \num{n_k}, y) \lif \num{l} =
  y)]$। तर्कशास्त्र और इस मान्यता का उपयोग करके कि
  $\Th{Q} \Proves !A_f(\num{n_0}, \dots, \num{n_k}, \num{m})$, हमें
  मिलता है कि $\Th{Q} \Proves \eq[\num{l}][\num{m}]$। दूसरी ओर,
  \olref[bre]{lem:q-proves-neq} के अनुसार
  $\Th{Q} \Proves \eq/[\num{l}][\num{m}]$। अतः $\Th{Q}$ असंगत है। किंतु
  यह असंभव है, क्योंकि $\Th{Q}$ मानक प्रतिरूप द्वारा संतुष्ट होता है
  (देखें \olref[int][def]{def:standard-model}), $\Sat{N}{\Th{Q}}$, और
  सुदृढ़ता प्रमेय के अनुसार संतुष्टनीय सिद्धांत सदैव संगत होते हैं
  (\tagrefs{prfAX/{fol:axd:sou:cor:consistency-soundness},
prfSC/{fol:seq:sou:cor:consistency-soundness},
prfND/{fol:ntd:sou:cor:consistency-soundness},
prfTab/{fol:tab:sou:cor:consistency-soundness}})।
\end{proof}

\begin{lem}
$\Th{Q}$ में निरूपणीय प्रत्येक फलन संगणनीय है।
\end{lem}

\begin{proof}
पहले यह सत्य क्यों है, इसका सहज विचार दें। $f$ को संगणित करने के लिए हम
यह करते हैं। अंकगणित की भाषा में सभी संभव !!{derivation}s~$\delta$ की सूची
बनाएँ। यह यांत्रिक रूप से किया जा सकता है। प्रत्येक के लिए जाँचें कि क्या
वह $!A_f(\num{n_0}, \dots, \num{n_k}, \num{m})$ रूप के किसी
!!a{formula} की !!a{derivation} है (अर्थात् \olref{lem:rep-q} का वह
!!{formula} जो~$\Th{Q}$ में $f$ का निरूपण करता है)। यदि हाँ, तो
\olref{lem:rep-q} के अनुसार $m = f(n_0, \dots, n_k)$, और हमें~$f$ का मान
मिल गया है। खोज समाप्त होती है, क्योंकि
$\Th{Q} \Proves !A_f(\num{n_0}, \dots, \num{n_k},
\num{f(n_0, \dots, n_k)})$; अतः अंततः हमें सही प्रकार का कोई~$\delta$ मिल
जाता है।

यह अभी पूरी तरह सटीक नहीं है, क्योंकि हमारी प्रक्रिया केवल संख्याओं के
स्थान पर !!{derivation}s और !!{formula}s पर काम करती है, और हमने ठीक-ठीक
यह नहीं बताया कि ``सभी संभव !!{derivation}s की सूची बनाना'' यांत्रिक रूप
से क्यों संभव है। किंतु जैसा हमने देखा है, पदों, !!{formula}s और
!!{derivation}s को ग्योडेल (G\"odel) संख्याओं से कूटबद्ध किया जा सकता है।
हमने संगणना का एक सटीक प्रतिरूप---सामान्य पुनरावर्ती फलन---भी प्रस्तुत
किया है। और हमने देखा है कि संबंध $\Prf[\Th{Q}](d,y)$ (जो तभी लागू होता है
जब $d$,~$\Th{Q}$ के अभिगृहीतों से ग्योडेल (G\"odel) संख्या~$y$ वाले
!!{formula} की !!a{derivation} की ग्योडेल (G\"odel) संख्या हो) (आदिम)
पुनरावर्ती है। हमें जिन अन्य आदिम पुनरावर्ती फलनों की आवश्यकता होगी, वे
$\fn{num}$ (\olref[art][trm]{prop:num-primrec}) और $\fn{Subst}$
(\olref[art][sub]{prop:subst-primrec}) हैं। इनसे~$f$ को न्यूनीकरण द्वारा
परिभाषित किया जा सकता है; अतः $f$ पुनरावर्ती है।

पहले परिभाषित करें
\begin{multline*}
  A(n_0, \dots, n_k, m) = \\
  \fn{Subst}(\fn{Subst}(\dots\fn{Subst}(\Gn{!A_f}, \fn{num}(n_0), \Gn{x_0}),\\ \dots),
  \fn{num}(n_k),  \Gn{x_k}), \fn{num}(m), \Gn{y}).
\end{multline*}
यह जटिल दिखता है, किंतु यह केवल फलन
$A(n_0, \dots, n_k, m) =
\Gn{!A_f(\num{n_0}, \dots, \num{n_k}, \num{m})}$ है।

अब संबंध~$R(n_0, \dots, n_k, s)$ पर विचार करें, जो तब लागू होता है जब
$(s)_0$,~$\Th{Q}$ से
$!A_f(\num{n_0}, \dots, \num{n_k}, \num{(s)_1})$ की !!a{derivation} की
ग्योडेल (G\"odel) संख्या हो:
\[
R(n_0, \dots, n_k, s) \quad\text{तभी और केवल तभी जब}\quad
\Prf[\Th{Q}]((s)_0, A(n_0, \dots, n_k, (s)_1))
\]
यदि हमें ऐसा~$s$ मिल जाए जिसके लिए $R(n_0, \dots, n_k, s)$ लागू हो, तो
हमें संख्याओं का ऐसा युग्म---$(s)_0$ और~$(s)_1$---मिल गया है जिसमें
$(s)_0$,~$A_f(\num{n_0}, \dots, \num{n_k}, (s)_1)$ की
!!a{derivation} की ग्योडेल (G\"odel) संख्या है। अतः~$s$ को खोजना अनौपचारिक
प्रमाण के $d$ और~$m$ युग्म को खोजने जैसा है। ऐसे~$s$ को ``खोजने वाला''
संगणनीय फलन नियमित न्यूनीकरण द्वारा परिभाषित किया जा सकता है। ध्यान दें कि
$R$ नियमित है: प्रत्येक $n_0$, \dots, $n_k$ के लिए
$\Th{Q} \Proves !A_f(\num{n_0}, \dots, \num{n_k},
\num{f(n_0, \dots, n_k)})$ की कोई !!a{derivation}~$\delta$ है; इसलिए
$s = \tuple{\Gn{\delta}, f(n_0, \dots, n_k)}$ के लिए
$R(n_0, \dots, n_k, s)$ लागू होता है। अतः हम~$f$ को इस प्रकार लिख सकते हैं:
\[
f(n_0,\dots,n_{k}) = (\umin{s}{R(n_0, \dots, n_k, s)})_1.
\]
\end{proof}

\end{document}
